\section{Black Holes, Entropy, and Hawking Radiation}

Black holes are where gravity, the quantum, and information collide, and they are the toughest test bench in physics. A good theory of gravity should say sensible things about them. The crystal does, recovering two of the most celebrated results about black holes, both of which were deep surprises when first found.

\subsection{Entropy lives on the surface, not the volume}

Entropy is a measure of hidden information, how many microscopic arrangements look the same from outside. For almost everything, entropy grows with volume: a box twice as big in every direction holds eight times as much hidden information. Black holes break this rule shockingly. A black hole's entropy grows with its surface area, not its volume. Double its size and the entropy goes up fourfold, not eightfold. The information seems to live on the horizon, the two-dimensional surface, not in the three-dimensional interior. This is the famous area law, and it is the seed of the idea that the world might be, in some sense, a hologram.

The crystal reproduces it. Take a region of the web and ask how much its inside is entangled with its outside, how much shared information crosses the boundary. On the crystal that shared information scales with the boundary, the number of notes crossing the surface, not with the volume inside. So the area law is not mysterious here. It is the natural behavior of information on a web, where what the inside and outside share is exactly what passes through the connecting notes at the surface. Area, not volume, because the surface is where the crossing notes are.

\subsection{Hawking radiation: black holes glow}

The second great surprise is Stephen Hawking's: black holes are not perfectly black. They glow, very faintly, with a thermal radiation, and slowly evaporate. Smaller holes are hotter. As a hole shrinks it gets hotter and glows faster, until it is gone.

The crystal recovers this too. Near a horizon, the quantum field is in a special twinned state across the boundary, and when you can only see one side, what you see is exactly a warm thermal glow, at the temperature Hawking predicted. The temperature comes out scaling the right way, hotter for smaller holes. And, importantly, when you track the information over the hole's whole life, it follows the curve that says information is not destroyed but gradually returned, a question that has tormented physics for fifty years. The crystal lands on the side of information preserved.

\subsection{Why black holes matter for the whole theory}

It would be easy to dismiss black holes as exotic edge cases. They are not. They are the sharpest place where a theory of discrete space, gravity, and the quantum must all hold together at once, and where most theories either go silent or give nonsense. That the crystal recovers the area law and Hawking's glow, from the same web and rule that gave us ordinary space and matter, is strong evidence that the pieces are fitting together rather than being patched.

There is also a satisfying unity here. The area law is about information on a web. Hawking's glow is about the quantum field near a horizon. The singularity-capping from the last chapter is about the smallest length. All three are the same crystal showing different faces, and all three come out right without separate assumptions.

\textbf{Takeaway:} the crystal reproduces the two great black-hole results: entropy that scales with surface area, not volume, because shared information lives on the crossing notes at the boundary, and Hawking's thermal glow, hotter for smaller holes, with information gradually returned rather than lost. These are the hardest joint test of discrete space, gravity, and the quantum, and the crystal passes.
